\chapter{La gravità dell'algoritmo e la caverna digitale}
\label{cap:p2c11}

\epigraph{«I social media sono bias cognitivi con le gambe: si muovono, si riproducono, si evolvono.»}{L'autore}

\section{Il mito dell'algoritmo rotto e la piramide di Wallace}

Nel capitolo precedente le macchine imparavano i nostri pregiudizi da archivi chiusi, dieci anni di curriculum o di cartelle cliniche. I sistemi di raccomandazione delle piattaforme fanno qualcosa di diverso. Il loro archivio si aggiorna in tempo reale, ed è fatto dei nostri clic, delle nostre pause davanti a un video, delle nostre condivisioni. Imparano da ciò che abbiamo appena scelto e decidono che cosa vedremo subito dopo. Qui la macchina non si limita a riprodurre il passato, lo restituisce su scala, a miliardi di persone, orientando ciò che finisce davanti ai loro occhi.

Il 10 dicembre 2025 l'Australia è diventata il primo Paese al mondo a vietare alle piattaforme social di permettere a chi ha meno di sedici anni di aprire o mantenere un proprio account. Nelle settimane successive le piattaforme hanno disattivato un gran numero di profili, molti ragazzi hanno trovato il modo di aggirare i controlli e in mezzo mondo si è discusso se seguire l'esempio. Comunque la si pensi, la decisione poggia su un'idea ormai diffusa, che questi ambienti non siano neutri e che per qualcuno possano essere dannosi.

Si dice spesso qualcosa di più, cioè che gli algoritmi dei social ci abbiano polarizzati, chiudendo ciascuno in una bolla. È una spiegazione intuitiva, persino rassicurante, perché individua un responsabile esterno a noi. Le prove migliori di cui disponiamo raccontano una storia più complicata.

Pochi anni prima del divieto australiano, un gruppo di ricercatori indipendenti aveva ottenuto da Meta qualcosa di raro, il permesso di modificare per tre mesi, con il consenso degli utenti, il funzionamento di Facebook e Instagram per decine di migliaia di persone durante le elezioni americane del 2020. In un esperimento il feed algoritmico fu sostituito da un semplice ordine cronologico; in un altro fu ridotta l'esposizione a fonti politicamente affini. In entrambi i casi cambiò ciò che le persone vedevano, e quanto tempo passavano sulle piattaforme, ma i loro atteggiamenti polarizzati, misurati con questionari, non cambiarono in modo misurabile.\footnote{Guess, A. M., e altri (2023). How do social media feed algorithms affect attitudes and behavior in an election campaign? \textit{Science}, 381(6656), 398-404; Nyhan, B., e altri (2023). Like-minded sources on Facebook are prevalent but not polarizing. \textit{Nature}, 620, 137-144. Gli autori stessi avvertono che tre mesi sono pochi per misurare effetti accumulati in anni.} Anche le prove sulle «bolle di filtraggio» sono più deboli di quanto si pensasse. Già nel 2015 uno studio sui dati di Facebook, condotto da ricercatori dell'azienda, aveva trovato che a limitare l'esposizione a notizie di orientamento opposto erano le scelte degli utenti, gli amici che avevano e i link su cui cliccavano, più dell'ordinamento algoritmico.\footnote{Bakshy, E., Messing, S., e Adamic, L. A. (2015). Exposure to ideologically diverse news and opinion on Facebook. \textit{Science}, 348(6239), 1130-1132.} Un'analisi della navigazione di decine di migliaia di utenti ha trovato che social e motori di ricerca aumentavano un po' la distanza ideologica tra le persone, ma anche la loro esposizione a idee contrarie, e che la maggior parte delle notizie si leggeva comunque sui siti delle grandi testate.\footnote{Flaxman, S., Goel, S., e Rao, J. M. (2016). Filter bubbles, echo chambers, and online news consumption. \textit{Public Opinion Quarterly}, 80(S1), 298-320.} E negli Stati Uniti la polarizzazione è cresciuta di più proprio tra gli anziani, cioè tra chi usa meno internet e i social.\footnote{Boxell, L., Gentzkow, M., e Shapiro, J. M. (2017). Greater Internet use is not associated with faster growth in political polarization among US demographic groups. \textit{Proceedings of the National Academy of Sciences}, 114(40), 10612-10617.}

Nessuno di questi studi assolve le piattaforme. Misurano periodi, Paesi e indicatori precisi, e lasciano aperte domande importanti sugli effetti accumulati, sui gruppi più esposti, su ciò che accade fuori dagli Stati Uniti. Ma escludono la versione più semplice della storia. Gli algoritmi contano, però non sono un telecomando che regola le menti di tutti allo stesso modo. Lavorano su tendenze che esistevano prima di loro e le amplificano più per alcuni che per altri. Il software fa quello per cui è stato ottimizzato; le vulnerabilità su cui fa leva sono nostre. È con questa doppia lente che guarderemo il resto del capitolo.

Per capire dove l'amplificazione trovi presa serve una mappa. Lo scrittore americano David Foster Wallace aveva intuito qualcosa sulla natura umana prima ancora che i social esistessero nella forma che conosciamo, una legge quasi fisica applicata alla sociologia: \textbf{le persone vengono unite in massa dai loro istinti più bassi, e separate dai loro interessi più elevati}. Immaginiamo allora la nostra specie come una gigantesca piramide, non di potere o di ricchezza, ma di complessità cognitiva.

Alla base si trova il più veloce dei due piloti, il Sistema 1\index[main]{Sistema 1} del cap.~\ref{cap:p1c2},\footnote{Kahneman, D. (2011). \textit{Pensieri lenti e veloci}. Mondadori, Milano.} rapido e inconsapevole, che reagisce alla paura, al desiderio, alla rabbia, alla curiosità verso il pericolo. A questo livello ci somigliamo molto più di quanto ci piaccia credere, per quanti anni abbiamo studiato e quanti libri abbiamo letto. Davanti a un incidente stradale o a una rissa, un premio Nobel e il concorrente di un reality mostrano reazioni di allarme simili, che coinvolgono strutture come l'amigdala.\footnote{LeDoux, J. (1996). \textit{The Emotional Brain}. Simon and Schuster, New York.} Qui cultura ed educazione pesano poco, e la paura e l'indignazione diventano un collante potente.

Al vertice troviamo qualcosa di completamente diverso, il Sistema 2 al massimo delle sue possibilità. È qui che vivono la poesia difficile, l'astrofisica, la filosofia, la musica da camera, la riflessione lenta e faticosa. Un territorio meraviglioso, ma con un difetto commerciale fatale, perché divide. La raffinatezza intellettuale ci frammenta in nicchie sempre più piccole. Trovare dieci milioni di persone che si emozionino simultaneamente per un trattato di epistemologia è quasi impossibile. La Figura~\ref{fig:piramide_wallace} mette in scena questa gerarchia.

\figurasola{images/piramide_wallace.png}{La piramide cognitiva dell'umanità: alla base gli istinti universali che ci uniscono, al vertice l'intelletto che ci distingue e ci separa. L'algoritmo punta sempre alla base, dove vive la massa.}{fig:piramide_wallace}

\section{La gravità digitale e la slot machine in tasca}

In questa sproporzione geometrica, tra la base larghissima e il vertice strettissimo, entra in gioco l'algoritmo. L'obiettivo primario di una grande piattaforma commerciale non è l'educazione, né il benessere degli utenti, né la qualità del dibattito pubblico. È catturare la nostra attenzione, trattenerla, monetizzarla. Il secondo livello del capitolo precedente, l'obiettivo, qui si vede a occhio nudo.

Per raggiungere decine di milioni di persone conviene puntare alla base della piramide, al minimo comune denominatore cognitivo. Le reazioni a un saggio filosofico sono lente, disomogenee, imprevedibili. Quelle alla paura, allo scandalo, all'indignazione morale sono rapide e diffuse. L'algoritmo non ci odia. Non ha un'agenda morale. Con la freddezza di un'equazione, premia ciò che ottiene reazioni misurabili. Esercita una sorta di forza di gravità digitale, che tende a tirare verso il basso. Ma la gravità non agisce allo stesso modo su tutti i corpi, e dipende dalla massa, cioè, fuor di metafora, da ciò che ciascuno cerca, dalle sue abitudini, dalle alternative che ha.

Come una corrente sottomarina, la sentiamo appena. Ogni singola scrollata sembra una nostra libera scelta, e tecnicamente lo è. Ma la corrente ha già reso più facile nuotare in una certa direzione.

\phantomsection\label{sec:p2c11-dopamina}
La gravità trova un alleato nel modo in cui i telefoni sono disegnati. Critici che venivano dall'industria, come l'ex esperto di etica del design di Google Tristan Harris, hanno paragonato il gesto di trascinare il dito per aggiornare il feed a quello di tirare la leva di una slot machine.\footnote{Harris, T. (2016). «How Technology Hijacks People's Minds: from a Magician and Google's Design Ethicist». \textit{Medium}, 18 maggio 2016.} Il principio è la ricompensa variabile. A volte, scorrendo, si trova qualcosa di interessante; a volte no. Dagli esperimenti di psicologia dell'apprendimento sappiamo che ricompense imprevedibili mantengono un comportamento più a lungo di quelle regolari. Sean Parker, ex presidente di Facebook, ha ammesso pubblicamente che la piattaforma era stata pensata per sfruttare una vulnerabilità della psicologia umana.

Qui occorre correggere un equivoco diffuso. La dopamina non è «la molecola del piacere». Le ricerche la collegano piuttosto all'anticipazione, alla motivazione e all'apprendimento dalle ricompense, in particolare da quelle inattese. È un'immagine utile per capire perché si possa continuare a scorrere senza sentirsi meglio, inseguendo qualcosa che è sempre alla scrollata successiva. Ma resta un'immagine. Il passaggio dal neurotrasmettitore al nostro comportamento davanti al feed non è stato misurato direttamente, e quando nel discorso pubblico si parla di «dosi di dopamina» si usa una metafora, non un dato.

\section{Le camere dell'eco e l'ombra sulla parete}
\label{sec:p2c11-camere-eco}

Le camere dell'eco, quelle stanze immaginarie dove le voci rimbalzano amplificandosi senza mai incontrare un contraddittorio, esistevano molto prima di internet. Ma erano limitate dalla geografia: il bar del paese, il circolo politico, la parrocchia. Con i social, chi lo desidera può costruirsene una su scala globale.\footnote{Sunstein, C. R. (2017). \textit{\#Republic: Divided Democracy in the Age of Social Media}. Princeton University Press, Princeton.} Per la maggior parte delle persone, come abbiamo visto, la stanza resta più aperta di quanto si tema; per una minoranza può chiudersi molto. E le aziende non sono spettatrici innocenti. Secondo i documenti interni resi pubblici da Frances Haugen, a Facebook si sapeva che certe scelte di ordinamento favorivano i contenuti più divisivi, e non sempre si scelse di intervenire.\footnote{Haugen, F. (2021). Testimonianza alla Commissione Commercio del Senato degli Stati Uniti.}

Ciò che la tecnologia ha costruito è la versione elettrica della più antica immagine filosofica dell'illusione. Quasi venticinque secoli fa, nel settimo libro della \textit{Repubblica}, Platone\index[main]{Platone} immaginò degli uomini incatenati fin dalla nascita in fondo a una caverna, costretti a fissare una parete.\footnote{Platone. \textit{La Repubblica}, libro VII. Numerose edizioni.} Alle loro spalle un fuoco proietta ombre di oggetti che essi non possono vedere direttamente. Per quei prigionieri le ombre sono l'unica realtà che conoscono. Hanno persino imparato a dare un nome alle ombre, a prevederne i movimenti, a competere su chi sappia riconoscerle più in fretta. Bastano poche modifiche per rendere quel mito contemporaneo. La caverna è lo schermo, il fuoco è la luce dei nostri dispositivi, le ombre sono i contenuti che l'algoritmo seleziona per noi in base a ciò che ci trattiene. La differenza è che le nostre catene sono fatte di abitudine anziché di ferro, e che la parete si aggiorna in tempo reale per non annoiarci mai.

\section{Il bias della disponibilità e il rabbit hole}

Su queste pareti luminose, l'informazione più visibile non è necessariamente la più accurata, ma spesso la più coinvolgente. Analizzando la diffusione di notizie su Twitter tra il 2006 e il 2017, uno studio ha trovato che quelle false arrivavano più lontano e più in fretta di quelle vere; gli autori collegano il fenomeno alla novità e alle reazioni emotive che suscitano, e osservano che a diffonderle erano soprattutto persone, non bot.\footnote{Vosoughi, S., Roy, D., e Aral, S. (2018). The spread of true and false news online. \textit{Science}, 359(6380), 1146-1151.} Il sistema amplifica, ma il primo clic è nostro.

C'è poi un meccanismo più subdolo. Il cervello valuta la probabilità di un evento in base alla facilità con cui riesce a richiamarlo alla memoria. È il bias della disponibilità\index[cognitive]{bias di disponibilità}, quello che ci fa temere gli squali dopo un film di paura ambientato in mare. Moltiplicato per un flusso continuo di drammi, scandali, disastri e crimini, può farci percepire un mondo più pericoloso e corrotto di quanto sia. Per molti indicatori di lungo periodo, dalla mortalità infantile alla violenza, viviamo in un'epoca più sicura di quasi tutte le precedenti, eppure ci sentiamo costantemente minacciati. La normalità fa cliccare poco. È come vivere in un mondo dove esistono soltanto incendi, e mai i pompieri che li spengono.

Dalla distorsione dei singoli contenuti si passa a quella dei percorsi. Alcune piattaforme video sono accusate di un'arte ulteriore, quella che in inglese si chiama \textit{rabbit hole\index[main]{rabbit hole}}, la tana del coniglio. Il racconto tipico è noto: si comincia con un innocuo video di gattini e, tre ore dopo, ci si ritrova tra i terrapiattisti. Le ricerche più accurate raccontano una storia meno cinematografica e più interessante. Analizzando la navigazione reale di migliaia di utenti, diversi studi hanno trovato che il consumo di contenuti estremi su YouTube è concentrato in una piccola minoranza di persone, che spesso arriva a quei canali da link esterni o abbonamenti, e che aveva già atteggiamenti affini; per la gran parte degli utenti le raccomandazioni non producono spirali di radicalizzazione.\footnote{Hosseinmardi, H., e altri (2021). Examining the consumption of radical content on YouTube. \textit{PNAS}, 118(32); Chen, A. Y., Nyhan, B., Reifler, J., Robertson, R. E., e Wilson, C. (2023). Subscriptions and external links help drive resentful users to alternative and extremist YouTube channels. \textit{Science Advances}, 9(35). Le piattaforme modificano spesso i loro algoritmi, e i risultati valgono per il periodo studiato.} \textbf{È il bias di conferma che incontra un sistema che lo asseconda}. L'algoritmo non trascina chiunque nella tana, ma rende molto facile scendere a chi ci sta già andando.

Questa discesa si innesta sul tribalismo. Secondo una stima molto citata, il numero di relazioni stabili che riusciamo a gestire si aggira intorno al centinaio e mezzo, e oggi navighiamo reti di milioni.\footnote{Dunbar, R. I. M. (1992). Neocortex size as a constraint on group size in primates. \textit{Journal of Human Evolution}, 22(6).} Il «noi contro loro» non riguarda più il villaggio vicino, ma può mettere milioni di persone contro altri milioni per ragioni sempre più astratte. Messaggi brevi che non lasciano spazio alle sfumature e strumenti che premiano la condivisione indignata non aiutano. Non a caso, le discussioni più feroci riguardano spesso argomenti di cui sappiamo pochissimo, perché è più facile odiare un'astrazione che una persona reale.

\section{La realtà che si sgretola}

Prima dei social esisteva qualcosa che potremmo chiamare realtà condivisa. Molti guardavano gli stessi pochi canali televisivi e leggevano gli stessi giornali. Si poteva dissentire sull'interpretazione, ma i fatti di base erano in buona parte comuni. Oggi il flusso di ciascuno può essere molto diverso da quello del vicino.\footnote{Pariser, E. (2011). \textit{The Filter Bubble: What the Internet Is Hiding from You}. Penguin Press, New York.} Gli studi citati all'inizio invitano a non esagerare. La bolla perfetta è rara, e la polarizzazione ha radici che precedono i social. Ma la frammentazione delle fonti è reale, e incontriamo persone intelligenti e colte convinte sinceramente di cose opposte sui fatti più elementari. Non è stupidità, è anche il risultato di ecosistemi informativi che si parlano sempre meno. Le due realtà, quella condivisa e quella digitale, sono messe a confronto nella Figura~\ref{fig:1000087949}.

\figurasola{images/1000087949.png}{Realtà condivise contro realtà digitali}{fig:1000087949}


Negli anni Venti del secolo scorso, il giornalista Walter Lippmann aveva descritto quello che chiamò \textit{pseudo ambiente}. Tra noi e il mondo reale, sosteneva, si frappone sempre una rappresentazione mentale, una mappa che non coincide mai del tutto con il territorio.\footnote{Lippmann, W. (1922). \textit{Public Opinion}. Harcourt, Brace and Company, New York.} Agiamo non sul mondo, ma su questa immagine del mondo. Lippmann pensava ai giornali e alla propaganda. Gli algoritmi rendono possibile uno pseudo ambiente costruito su misura per ciascuno. Un secolo fa il rischio era che molti condividessero la stessa mappa distorta; oggi è che ne condividiamo sempre meno.

Ed è qui che la questione diventa politica. Il filosofo tedesco Jürgen Habermas ha dedicato la vita a studiare la sfera pubblica e quello che chiamava agire comunicativo, lo scambio in cui cittadini liberi deliberano insieme lasciando che sia la forza del miglior argomento, e non quella del più forte, a orientare le decisioni comuni.\footnote{Habermas, J. (1981). \textit{Theorie des kommunikativen Handelns}. Suhrkamp, Francoforte.} Una deliberazione del genere ha una condizione elementare, che diamo per scontata solo finché esiste. I partecipanti devono condividere un minimo di fatti su cui ragionare. Si può discutere all'infinito su cosa sia giusto fare, ma non si può deliberare se non si concorda nemmeno su cosa stia accadendo. Senza un mondo comune, la conversazione democratica non si svolge male, semplicemente fatica a cominciare.

\section{Il ciclo: dove agisce la persona, dove agisce il sistema}
\label{sec:p2c11-ciclo}

Vale la pena fermarsi un momento per mettere in ordine i pezzi. La Figura~\ref{fig:ciclo_bias} mostra il percorso di un'informazione, dal momento in cui ci viene mostrata a quello in cui la nostra reazione torna a influenzare ciò che vedremo dopo. Alcuni passaggi avvengono nella nostra testa, altri nei sistemi delle piattaforme. Distinguerli serve a capire dove si può intervenire, sulle proprie abitudini, sulla progettazione dei sistemi, o su entrambe.

\begin{figure}[htbp]
\centering
\begin{tikzpicture}[
  node distance=5mm and 4mm,
  box/.style={draw, rounded corners=2pt, align=center, font=\sffamily\footnotesize, minimum width=23mm, minimum height=13mm, text width=22mm, fill=white},
  arr/.style={-{Stealth[length=2.2mm]}, thick},
  lab/.style={font=\sffamily\scriptsize\bfseries, text=moderngray}
]
% persona (riga in basso)
\node[box] (att) {\textbf{Attenzione}\\ notiamo ciò che è nuovo, emotivo, familiare};
\node[box, right=of att] (int) {\textbf{Interpretazione}\\ conferma, alone, disponibilità};
\node[box, right=of int] (dec) {\textbf{Decisione}\\ clic, condivido, credo, scelgo};
% sistema (riga in alto)
\node[box, above=11mm of att, fill=aicolor!8] (inf) {\textbf{Informazione}\\ selezione e ordine dei contenuti};
\node[box, above=11mm of dec, fill=aicolor!8] (rin) {\textbf{Rinforzo}\\ il sistema registra la reazione e aggiorna le previsioni};
% frecce
\draw[arr] (inf) -- (att);
\draw[arr] (att) -- (int);
\draw[arr] (int) -- (dec);
\draw[arr] (dec) -- (rin);
\draw[arr] (rin) -- node[above, font=\sffamily\scriptsize, align=center]{ciò che vedremo dopo} (inf);
% fasce
\begin{scope}[on background layer]
\node[fit=(inf)(rin), fill=aicolor!4, draw=aicolor!40, rounded corners=4pt, inner sep=7pt, label={[lab]above:SISTEMA DIGITALE (piattaforma, algoritmo)}] {};
\node[fit=(att)(int)(dec), fill=lightgray!60, draw=moderngray!40, rounded corners=4pt, inner sep=7pt, label={[lab]below:ESSERE UMANO (bias cognitivi)}] {};
\end{scope}
\end{tikzpicture}
\caption{Il ciclo di interazione tra bias umani e sistemi digitali. In grigio i passaggi che avvengono nella mente, in azzurro quelli che avvengono nel sistema. Le frecce indicano influenze possibili, non causalità automatiche: quanto ciascun passaggio pesi davvero varia molto da persona a persona e da piattaforma a piattaforma, come mostrano gli studi citati in questo capitolo.}
\label{fig:ciclo_bias}
\end{figure}

Nessun anello del ciclo è sufficiente da solo. Senza la nostra attenzione selettiva, il sistema non avrebbe nulla da rinforzare; senza il rinforzo, le nostre tendenze resterebbero confinate alle nostre scelte. È nella somma che il ciclo può stringersi. Ed è nella somma che si vede la domanda di questo libro in azione. Una mente che distorce la realtà produce clic distorti, i clic diventano dati, i dati addestrano il sistema, e il sistema decide quale realtà mostrerà domani a quella stessa mente. Per questo le risposte efficaci, come vedremo nella Parte IV, lavorano su entrambi i lati.

\section{Il pensiero come unica via d'uscita}

Come si esce da questa macchina? Non possiamo disinstallare i nostri bias. Ma c'è qualcosa di unico nella mente umana, la capacità di osservare i propri processi mentali mentre accadono. Si chiama \textbf{metacognizione}, letteralmente pensare al proprio pensiero. È insieme la nostra condanna e la nostra salvezza. Condanna, perché ci rende consapevoli di quanto siamo inadeguati al mondo che abbiamo costruito. Salvezza, perché una volta che si vedono i fili del burattinaio, la magia perde una parte del suo potere.

Vale la pena ascoltare una voce che ha attraversato il secolo più buio. La filosofa Hannah Arendt\index[main]{Arendt, Hannah}, riflettendo sulle catastrofi politiche del Novecento, giunse a una conclusione sorprendente. Alla radice del male su larga scala non c'è quasi mai un'intelligenza diabolica, ma molto più spesso una semplice, diffusa incapacità di pensare.\footnote{Arendt, H. (1978). \textit{The Life of the Mind}. Harcourt Brace Jovanovich, New York.} Non l'assenza di intelligenza, si badi, ma la rinuncia a fermarsi, a interrogarsi, soprattutto a guardare il mondo dal punto di vista di un altro. È la facoltà che una caverna su misura rischia di atrofizzare, perché ci circonda di voci che ci somigliano e rende l'altro invisibile, o peggio, una caricatura da disprezzare.

Cosa farne, in concreto? Come in una dieta equilibrata, anche nell'informazione la varietà conta. Seguire deliberatamente fonti che ci sfidano è scomodo come l'esercizio fisico, e altrettanto necessario. Non occorre per forza eliminare i social, anche se per qualcuno può essere la scelta migliore. Serve piuttosto una sorta di igiene digitale, con zone franche senza dispositivi: la prima ora del mattino, i pasti, almeno un giorno alla settimana. Periodicamente conviene potare le proprie fonti. E soprattutto, prima di condividere un contenuto che provoca una reazione emotiva violenta, fermiamoci un istante e chiediamoci quale dei due piloti ha preso il volante. È informazione, o è indignazione confezionata per generare un clic?

C'è un'ultima via d'uscita, la più antica di tutte, e ce la indica ancora il mito della caverna. Per sottrarsi alla gravità dell'algoritmo occorre un atto che va contro ogni istinto sociale, accettare la solitudine intellettuale come un prezzo non soltanto necessario, ma nobile. Spegnere il rumore di fondo, rifiutare l'indignazione preconfezionata, coltivare le proprie altezze cognitive può farci sentire terribilmente isolati. Ma è in larga parte un'illusione ottica. Nel mito, chi spezza le catene e risale fuori dalla caverna soffre, accecato dal sole, e quando torna a raccontarlo agli altri viene deriso. Eppure ha visto. E se avremo la pazienza di sopportare questa temporanea solitudine, cominceremo a incontrare i nostri veri simili, non quelli che condividono i nostri riflessi, ma quelli che condividono le nostre domande, i nostri dubbi, le nostre passioni più autentiche. È l'unica forma di connessione che l'algoritmo non riuscirà mai a simulare.

Il feed, però, resta una parete su cui scorrono cose. Lo percepiamo come un flusso, non come qualcuno. Gli assistenti conversazionali compiono un passo in più, perché non si limitano a scegliere che cosa vediamo ma ci rispondono, in prima persona, nella nostra lingua. E davanti a qualcosa che parla come una persona entra in gioco un'altra delle nostre tendenze, quella di vedere una mente dove c'è soltanto un meccanismo. Se l'algoritmo del feed sfrutta ciò che ci attira, l'interlocutore artificiale può sfruttare ciò di cui abbiamo bisogno, cioè essere ascoltati.
